% Cote 119: offprints, reports, courses, CV and letters (typed and printed part): archive cote 119, site page 220 (PDF page 221).
% typescript. Draft compared: none.
% AI-transcribed by Claude Opus 5.5 (claude-opus-5-5, adaptive thinking, effort high), Message Batches API, batch r119_full (msgbatch_01DwZNm5Y7YsKBWyAN4sFs8w), 2026-11-13; not yet checked by eye.
\scanpage{119/220}
\mspage{. 7 .}

finitude (de Frisch et Siu) pour les images directes
supérieures d'un faisceau cohérent par une immersion
ouverte, utilisant la profondeur du faisceau en les
points du complémentaire, sont inspirés de théorèmes
analogues en géométrie algébrique [16, SGA 2] ; remarques
analogues pour des théorèmes sur la cohomologie à
supports compacts des faisceaux algébriques cohérents,
complétés par un théorème d'existence, et leur inter-%
prétation en termes de théorèmes du type de Lefschetz
pour la cohomologie cohérente (la version algébrique
fait partie de la thèse de Mme Michèle Raynaud (en cours
de publication), et la version analytique est due à
Trautmann [*]). Parlant en termes de grands thèmes de
recherche plutôt qu'en termes de résultats techniques
particuliers, je pense qu'outre les thèmes déjà nommés,
les thèmes suivants ont été directement suscités ou
tout au moins influencés par les idées que j'avais dé-%
veloppées en géométrie algébrique :

a) \uline{Techniques de construction d'espaces analytiques},
aboutissant aussi bien à des espaces ``modulaires''
``globaux'' comme les espaces modulaires de Picard, pour
certains espaces analytiques compacts comme dans [11]
(le cas général ne semble pas encore traité), qu'à des
espaces modulaires ``locaux'' de déformation d'une structure
analytique complexe donnée, ou au modèle de la Géométrie
Formelle (``th. d'existence des modules formels'',
cf. [15 bis, Exp. n° 195]). Dans certains cas, les
énoncés obtenus en géométrie algébrique sont directement
applicables (cf. M. Hakim [*]), dans d'autres de nouvelles
difficultés surgissent, pas toujours surmontées à
l'heure actuelle. Parmi les travaux définitifs dans
ce sens, on peut citer la thèse de A. Douady [*].

\margin{220}
